\xchapter{平面几何问题求解系统}{Plane geometry problem solving system}

% 前文构建了面向长步骤平面几何推理的多模态数据基准，为复杂几何任务提供了统一的评测基础；同时还提出了面向长链条平面几何推理的两阶段增强方法，为几何问题自动求解提供了有效的技术支撑。
为了进一步验证前述研究成果在实际应用场景中的可行性，并推动相关方法由离线实验走向系统落地，本文在前述工作的基础上开发了平面几何问题求解系统，对多模态输入、求解模型、数据集展示与结果可视化流程进行了统一集成。本章围绕系统的需求分析、系统设计与实现、系统展示与测试等内容展开说明，重点介绍几何求解模块的工程化实现过程。

\xsection{系统需求分析}{System Requirement Analysis}

\xsubsection{功能性需求分析}{Functional Requirement Analysis}

平面几何问题求解系统主要面向中学数学教学辅助、几何推理算法验证以及多模型实验对比等应用场景。与仅输出最终答案的传统数学工具相比，本系统不仅关注求解结果的正确性，也强调题目输入方式的适配能力、不同求解模型的统一管理能力，以及求解过程和结果的可视化展示。因此，在功能设计上，系统需要围绕题目输入、模型调用、结果返回和异常反馈等方面展开。

在题目输入方面，系统应支持用户输入平面几何题目的文字描述，并上传与题目对应的几何图像，以贴近实际试题中“题干+配图”的常见组织形式。为保证输入过程的完整性与便利性，前端界面需要提供清晰的文本输入区域、图像上传入口以及图片预览功能，后端则负责接收并处理文本与图像数据，为后续多模态求解提供基础输入。

在模型使用方面，系统应能够在统一框架下支持多种几何求解模型的配置与调用，使用户可以根据需要选择不同模型进行实验验证或结果比较。前端应提供便捷的模型勾选与切换方式，后端应支持不同模型运行环境的独立管理与统一调度，从而实现同一题目在不同求解方法下的横向对比，便于观察不同模型在求解效果、响应速度和稳定性等方面的差异。

在求解流程方面，系统应能够在用户提交题目后，将文本、图像和模型配置一并发送至后端服务，并由后端对所选模型进行并行调度。考虑到不同模型在运行时间、资源消耗以及外部依赖上的差异，系统不必等待全部模型完成后再统一返回结果，而是允许已完成的模型优先返回并展示其输出内容。这样可以有效降低用户等待时间，提升系统交互过程中的响应效率与使用体验。

在异常处理方面，由于题目输入、模型推理过程和网络通信都可能出现不确定情况，系统还应具备较为完善的状态识别与反馈能力。题目输入情况可能存在输入为空、图像缺失或者模型未选择，而网络异常可能导致无法正常接收和返回数据，同时服务端的后端依赖环境可能会出现错误，因此系统应具备较为完善地异常处理能力。对于不同的错误，系统应给出明确的状态提示和错误信息；对于已经成功完成求解的任务，则应及时展示结果。通过对成功任务和失败任务进行区分处理，系统能够在保证交互连续性的同时，提高整体可用性与稳定性。

\xsubsection{非功能性需求分析}{Non-Functional Requirement Analysis}

除功能性需求之外，平面几何问题求解系统还需满足教学辅助、算法验证和实验对比等场景下的若干非功能性要求，以保证系统在实际运行过程中具备良好的可用性、响应效率、稳定性以及后续演进能力。非功能性需求虽然不直接规定系统应提供哪些具体功能，但会直接影响系统的使用体验、服务质量和持续维护效果，因此同样是系统设计中的重要内容。

在易用性方面，系统界面应尽量符合用户的操作习惯，保证整体流程清晰、交互方式简洁，使用户能够较为直观地完成题目输入、模型选择和结果查看等操作。对于输入不规范、参数缺失或求解失败等情况，系统还应提供清晰明确的提示信息，以降低使用门槛并提升交互流畅性。

在响应性能方面，系统应具备较好的处理效率，能够在单模型或多模型求解场景下保持较为合理的响应时间。尤其在并发调用多个求解模型时，系统应尽量减少用户等待成本，保证结果展示过程具有较好的实时性和连贯性，从而提升整体交互体验。

在稳定性方面，系统应能够在多次重复调用、不同规模题目输入以及多模型同时运行等条件下保持正常服务状态。对于局部模块运行异常、外部依赖失败或个别求解任务中断等情况，系统应具备一定的容错能力，避免单点异常对整体服务造成明显影响，以保证系统运行的连续性与可靠性。

在可扩展性方面，系统应具备良好的模块扩展能力，以适应平面几何求解算法持续迭代的发展需求。未来新增模型或相关功能时，系统应能够以较低成本完成接入和集成，并尽量减少对现有模块的影响，从而为后续算法更新和功能完善预留空间。

在可维护性方面，系统整体结构应保持较好的模块化和低耦合特征，使界面层、业务逻辑层、模型服务层以及数据管理层之间具有清晰的职责划分。这样不仅有利于后续调试与故障定位，也有助于系统升级、功能优化和长期维护工作的开展。
\xsection{系统设计与实现}{System Design and Implementation}

\xsubsection{系统总体架构}{Overall System Architecture}

为满足多模态几何题在线求解与多模型对比分析的业务需求，系统在整体架构上采用前后端分离的分层式设计思路，将界面展示、交互通信、应用服务、模型执行以及数据存储进行解耦。该设计不仅有助于提升页面响应效率，也便于后端独立部署不同求解算法与运行环境。系统总体架构图如~\ref{fig:system_arch}所示。

\begin{figure}[htbp]
    \centering
    \includegraphics[width=0.95\textwidth]{c5_system_arch.pdf}
    \caption{系统总体架构}
    \label{fig:system_arch}
\end{figure}

系统整体可划分为前端展示层、交互控制层、后端服务层与数据支撑层四个部分。前端展示层负责用户输入、页面布局、结果渲染；交互控制层负责基于 HTTP 和 JSON 协议完成异步请求发送、参数封装与返回结果解析；后端服务层负责路由管理、输入校验、任务分发与结果聚合，同时包含模型执行、负责调用具体求解器并返回结构化结果；数据支撑层则承担题目记录、运行日志、数据集元信息持久化及热点缓存管理等功能，可以为用户多次运行提供缓存支持。

在整体数据流上，用户首先在前端页面中输入题目文本、上传几何图像并选择一个或多个求解模型；随后前端通过异步请求将输入数据发送至后端接口；后端完成参数合法性检查后，根据模型配置将任务分发到不同求解器；各求解器在独立运行环境中执行推理任务，并将结果或错误信息返回后端；后端服务层对已完成结果进行统一封装后，分批发送至前端；前端根据返回字段实时更新各模型对应的展示区域。该流程实现了“输入--调度--求解--展示”的完整链路，并通过并行返回机制增强了系统在多模型场景下的交互效率。

\xsubsection{系统核心模块}{Core Module}

围绕平面几何问题求解这一核心任务，系统在功能实现上主要包括输入管理、模型选择与调度、求解执行、以及结果展示四个部分。各部分在功能上相互衔接，共同构成从数据输入、模型调度、算法求解和结果展示的完整处理流程。

输入管理模块主要负责接收用户提交的题目文本、几何图像以及可选的数据集样例，并完成基础的格式检查、文件接收与请求封装。在图文联合输入场景下，该模块需要对文本信息与图像信息进行统一组织，使其能够以标准化形式传递至后端服务，从而为后续求解流程提供一致的输入基础。

模型选择与调度模块是系统运行过程中的核心控制单元。前端允许用户在同一界面中选择一个或多个求解模型，后端则依据用户的选择生成相应的调度任务。为支持不同求解器的统一管理，系统为各类模型维护了规范化的配置信息，包括模型标识、运行环境、调用方式以及输入输出规则等内容。基于这一机制，系统既能够保证现有模型的稳定调用，也能够为后续新增算法的接入提供便利。

求解执行模块负责调用具体的几何求解算法，并完成实际的推理计算。由于不同求解器对运行环境和依赖库的要求并不一致，系统在实现中采用了基于虚拟环境的隔离式部署方案，即为不同类型的求解器配置独立的 Conda 环境，并通过脚本接口完成模型启动、参数传递与结果回收。后端调度器只需依据预设配置调用相应脚本，便可以完成不同算法之间的切换与适配，从而有效降低环境冲突带来的影响，并提升系统的可扩展性与维护效率。

结果展示模块主要负责对后端返回的求解内容进行统一整理和可视化呈现。考虑到不同模型在输出形式上的差异，系统首先在后端对返回结果进行结构化处理，将其整理为统一的数据格式，主要包含模型名称、执行状态、求解答案、关键推理过程、运行耗时以及异常信息等内容。前端在接收到这些结果后，为每个模型生成相应的展示区域，从而实现多模型结果的独立呈现与横向比较。

\subsubsection{前后端通信}

系统前后端之间采用基于 HTTP 的异步方式进行数据交换。前端通过 Axios 向后端 FastAPI 接口发送请求，请求中包含题目文本、几何图像以及用户所选择的模型信息。后端接收到请求后，首先完成必要的参数校验与文件检查，随后将合法请求提交至调度模块，由其进一步分发给对应的求解器执行。

由于系统支持多个模型同时参与求解，前后端通信并未采用传统的阻塞式返回方式，而是按照异步处理思路设计为分阶段响应机制。具体而言，后端在各模型执行过程中持续维护任务状态，并将已经完成的结果优先返回给前端；前端在接收到局部结果后，仅更新对应模型的展示区域，而无需等待全部任务完成。对于仍在执行的模型，页面保持运行中的状态提示；若某一模型在执行过程中出现异常，则在对应区域显示失败信息，而不会影响其他模型已完成结果的查看。

这种通信方式能够更好地适应多模型并行求解场景。一方面，运行较快的模型可以优先输出结果，从而减少用户的感知等待时间；另一方面，即使某些模型执行较慢或出现异常，也不会阻塞整个页面的结果展示，因此能够在一定程度上提升系统的交互效率和运行稳定性。

\subsubsection{后端服务与模型部署}

系统后端服务主要基于 FastAPI 技术构建，用于承载几何实际求解相关的服务逻辑。其主要职责包括接收前端请求后完成输入校验，转入组织不同算法进行求解任务，并且调用具体模型，最后将执行结果以统一格式返回前端。相较于传统后端框架，FastAPI 在异步处理和接口组织方面具有较好的适配性，因此能够较好地满足模型调用型系统的实现需求。

在模型部署层面，系统主要依托 Python 生态完成各类几何求解器的封装与运行，并结合 PyTorch、Transformers 等深度学习框架支持多模态模型与大语言模型的推理过程。由于不同求解算法在底层依赖、Python 解释器版本以及具体调用方式上往往存在差异，因此系统并未将所有算法统一部署在同一环境中，而是采用相互隔离的部署策略。具体而言，系统为不同类型的求解器配置独立的 Conda 环境，并通过脚本接口完成模型启动、输入传递与结果回收。后端控制器在接收到任务后，只需依据预设配置调用对应环境中的执行脚本，即可完成相应算法的求解过程。

这一实现方式使系统在模型管理方面具有较好的灵活性。对于已经接入的算法，只需维护其运行环境和脚本入口，便能够保持稳定调用；对于新增算法，则可以在不改变整体业务流程的前提下，通过补充环境配置和适配脚本完成快速接入。

在异常处理方面，后端会对模型执行过程中可能出现的异常情况进行统一捕获，例如执行超时、外部依赖调用失败、脚本运行异常或返回结果格式不符合要求等。出现此类情况时，系统会将错误信息整理为统一响应格式后返回前端。通过这一机制，即使个别求解器无法正常完成任务，其余模型仍可继续运行并返回结果，从而避免局部故障对整体求解流程造成过大影响。

\subsubsection{数据存储与缓存}

在数据存储层，系统采用 MySQL 作为主要的存储数据库，用于保存系统运行过程中产生的结构化信息。这些数据主要包括用户提交的题目内容、图像文件路径、求解任务记录以及结果相关的元数据等。由于 MySQL 具有较为成熟的事务支持机制和稳定的数据管理能力，因此适合作为系统的持久化存储组件。

为提高系统在重复访问和连续请求场景下的响应效率，系统同时引入 Redis 作为缓存支持。与数据库主要承担长期保存功能不同，Redis 更适合用于维护短时状态和热点数据。例如，在相同题目被重复提交时，系统可以优先检查缓存中是否已存在对应结果；若结果可直接复用，则无需再次触发完整求解流程，从而降低后端计算开销。在多模型并行求解过程中，Redis 还可以辅助保存部分中间状态，使前后端之间的异步交互更加及时和平滑。

除服务于求解流程外，数据存储模块还承担了基础的数据资源管理任务。对于系统内置的数据集样例以及用户本地上传的数据文件，数据库会保存其必要的索引信息与映射关系，以支持前端的浏览与登记功能。由于本章重点在于几何求解模块的实现，因此该部分仅保留基本的数据支撑能力，其设计保持相对简化。

\xsection{系统展示与测试}{System Display and Testing}

\subsubsection{页面展示}

为了便于用户在统一平台中完成几何任务的使用与切换，系统在界面层面采用了“左侧导航栏 + 右侧功能区”的整体布局。首页作为整个几何项目的统一入口，主要用于展示系统标识及各功能模块之间的关系，如图~\ref{fig:system_homepage}所示。页面左侧为纵向导航栏，顶部展示“GeoSeek 几何求索”的名称与图标，下方依次设置“首页”“几何生成”“几何求解”和“数据集展示”等标签页，用户可通过导航栏在不同页面之间切换。右侧为首页主体区域，上方居中展示 GeoSeek 标识及几何主题图形，下方则以模块入口的形式给出“几何生成”“几何求解”和“数据集展示”三个功能板块。整体来看，首页布局较为简洁，功能划分也较为明确，能够直观呈现系统的基本结构。

\begin{figure}[htbp]
    \centering
    \includegraphics[width=0.95\textwidth]{c5_first_page.png}
    \caption{系统首页界面}
    \label{fig:system_homepage}
\end{figure}

首页的作用不仅提供页面跳转入口，也能从整体上组织系统内部的功能关系。其中，“几何生成”模块主要用于几何题目及图像样例的构造，可支持随机生成，也可依据预设范式生成相应样例，其在系统中承担了样例构造与任务准备的辅助作用。“几何求解”模块对应本章核心开发模块，是平面几何问题求解方法的主要承载页面。“数据集展示”模块则主要用于数据资源的浏览与查看，在系统中起到展示分析的作用。通过首页的统一组织，系统各模块之间形成了较清晰的层次关系，也使几何求解功能能够被放置在完整项目框架下加以理解。

在此基础上，平面几何问题求解页面采用“左侧输入区 + 右侧结果区”的双栏布局，在统一界面风格下完成题目输入、模型配置以及结果展示。页面左侧保留系统导航栏，主体区域顶部显示“几何问题求解器”标题，页面的主要交互过程围绕“输入--求解--展示”这一流程展开。

在页面初始状态下，左侧输入区主要包括文本输入、图片输入和模型选择三个部分，如图~\ref{fig:solver_page_input}所示。其中文本输入区域用于填写几何题题面，支持直接录入或粘贴题目描述；图片输入区域用于上传与题目对应的几何图像，并在页面中给出预览；模型选择区域提供四类求解模型的勾选入口，包括树状搜索求解、演绎可解释求解、大模型求解和神经符号求解。用户完成题面输入、图像上传与模型勾选后，可通过“开始求解”按钮向后端发起求解请求。

\begin{figure}[htbp]
    \centering
    \includegraphics[width=0.95\textwidth]{c5_solver_page_input.png}
    \caption{求解页面初始界面}
    \label{fig:solver_page_input}
\end{figure}

从页面组织方式来看，左侧输入区将题目文本、几何图像和求解模型集中放置，便于用户在一次交互中完成输入与配置；右侧结果区则预先为所选模型保留展示位置。在图~\ref{fig:solver_page_input}中，用户已勾选四种求解模型，因此右侧对应显示四个待返回的结果卡片，分别对应树状搜索求解、演绎可解释求解、大模型求解和神经符号求解。这样的界面安排能够使用户在提交前明确本次任务的求解范围，同时也为后续的结果对比提供统一的展示基础。

当用户点击“开始求解”后，系统在右侧结果区逐步呈现各模型的求解输出，如图~\ref{fig:solver_page_result}所示。结果区以模型为单位组织内容，不同模型的展示形式与其求解特点相对应。树状搜索求解结果更强调形式化解析与结构化推理，例如点、线、圆及其关系的形式表示，并结合图像中的关键元素高亮或局部识别结果，以体现其对图形结构的搜索能力；演绎可解释求解结果则以文字证明链为主，按照步骤给出条件、定理调用与推导过程，突出推理过程的可读性；大模型求解结果主要采用自然语言组织答案与分析，更接近日常题解的表述方式；神经符号求解结果则偏向解析几何式的计算过程，体现符号约束与数值推导相结合的求解路径。

\begin{figure}[htbp]
    \centering
    \includegraphics[width=0.95\textwidth]{c5_solver_page_result.png}
    \caption{求解结果展示界面}
    \label{fig:solver_page_result}
\end{figure}

由图~\ref{fig:solver_page_result}可以看出，结果区采用纵向卡片的方式组织多模型输出。当前已展开的树状搜索求解结果面板中，页面展示了题目的求解步骤、图形解析结果以及附加属性信息；其下方的演绎可解释求解结果继续以独立面板形式呈现，并在标题区域标识当前状态。这样的设计既保证了不同模型结果的独立性，也便于用户在同一页面内进行直接比较，从而更清楚地观察不同方法在同一道几何题上的表现差异。

此外，为了增强结果的可比较性，系统在各模型结果末尾统一附加四元组标签信息，即难度、求解步数、求解时间和涉及定理。该部分内容能够在较高层面对当前结果作出概括，一方面反映题目与推理过程的复杂程度，另一方面也提供了运行效率和知识调用方面的补充信息，从而为模型分析和案例比较提供辅助依据。总体而言，页面展示模块实现了从首页入口到结果呈现的完整交互过程，能够较好支撑平面几何问题的多模型求解与结果对比。

\subsubsection{功能测试}

为验证系统原型的可用性与稳定性，本文围绕平面几何问题求解流程中的关键环节开展了功能测试，重点考察系统在输入处理、模型调度、结果返回、页面展示以及异常处理等方面的表现。测试内容尽量覆盖用户实际使用过程中的主要操作路径，以检验系统是否能够完成从题目提交到结果展示的完整流程。测试项目与结果如表~\ref{tab:function_test}所示。

\begin{table*}[htbp]
    \small
    \centering
    \caption{系统主要功能测试结果}
    \label{tab:function_test}
    \begin{tabularx}{\textwidth}{
        >{\raggedright\arraybackslash}p{7em}
        >{\raggedright\arraybackslash}X
        >{\raggedright\arraybackslash}X
        c
        }
        \toprule
        测试项     & 测试输入                      & 预期结果                   & 实际结果 \\
        \midrule
        文本输入    & 输入合法几何题面                  & 题面信息能够被正确提交并传递至后端      & 通过   \\
        图像上传    & 上传 jpg/png 图像             & 图像能够完成上传并在页面中正常预览      & 通过   \\
        多模型选择   & 同时勾选 2--4 个求解模型           & 记录所选模型并向后端发起对应任务       & 通过   \\
        多模型并行求解 & 同一题目同时调用多种求解器             & 后端并行执行，各模型结果独立返回       & 通过   \\
        结果逐步渲染  & 部分模型先完成、部分模型仍在运行          & 前端按完成顺序更新结果区域，无需等待全部结束 & 通过   \\
        模型超时/失败 & 模拟单个模型超时或脚本执行失败           & 面板返回错误状态，其他模型继续正常运行    & 通过   \\
        异构结果展示  & 不同模型返回形式化解析、证明链、自然语言与计算过程 & 前端能够分类展示不同风格的求解结果      & 通过   \\
        统一标签展示  & 不同模型返回难度、步数、时间、涉及定理等信息    & 页面统一展示四元组标签，便于横向对比     & 通过   \\
        \bottomrule
    \end{tabularx}
\end{table*}

从测试结果可以看出，系统在文本输入、图像上传、模型选择、求解触发和结果展示等基础功能上均能稳定运行，各项测试结果均符合预期。对于平面几何问题求解这一任务，系统已经能够完成图文联合输入、多模型统一调度和异构结果展示的基本流程，说明其具备较好的工程可用性和任务适配能力。尤其是在题目提交与模型配置环节，系统能够较为准确地接收用户输入，并将相关信息传递给后端服务，为后续求解提供稳定的输入基础。与此同时，前端也能够根据模型返回内容及时更新页面，使用户可以直接查看不同模型的求解输出和展示形式。整体来看，系统已经具备原型平台应有的基本功能闭环。

在多模型并行求解和异常处理测试中，系统也表现出较好的稳定性。当多个求解器同时运行时，结果能够按照完成顺序逐步返回，前端无需等待全部任务结束即可开始展示内容，这在一定程度上提升了交互效率。当单个模型出现超时或调用失败时，系统能够正确反馈错误状态，且不会影响其他模型继续执行和展示结果。由此可以看出，系统在任务调度和异常隔离方面具有较好的可靠性，能够保证整体流程的连续性。结合表~\ref{tab:function_test}中的测试结果，可以认为该系统已经能够较好支撑平面几何问题的求解、对比与展示需求，达到了原型系统设计的基本目标。

\xsection{本章小结}{Summary}
本章介绍了平面几何问题求解系统的需求分析、系统设计、功能实现、页面展示与功能测试。围绕图文输入、模型选择、并行求解和结果展示等核心功能，完成了系统总体架构、功能模块、前后端通信、模型调度及数据存储等内容的设计与实现。结合页面展示与测试结果可以看出，所构建的系统能够支持平面几何题的输入、求解与结果输出，具备多模型统一接入、并行返回、异构结果展示和异常隔离等基本能力，为前文相关研究内容的系统集成与后续应用提供了基础。